\let\negmedspace\undefined
\let\negthickspace\undefined
\documentclass[journal]{IEEEtran}
\usepackage[a5paper, margin=10mm, onecolumn]{geometry}
%\usepackage{lmodern} % Ensure lmodern is loaded for pdflatex
\usepackage{tfrupee} % Include tfrupee package

\setlength{\headheight}{1cm} % Set the height of the header box
\setlength{\headsep}{0mm}     % Set the distance between the header box and the top of the text

\usepackage{gvv-book}
\usepackage{gvv}
\usepackage{cite}
\usepackage{physics}
\usepackage{amsmath,amssymb,amsfonts,amsthm}
\usepackage{algorithmic}
\usepackage{graphicx}
\usepackage{textcomp}
\usepackage{xcolor}
\usepackage{txfonts}
\usepackage{listings}
\usepackage{enumitem}
\usepackage{mathtools}
\usepackage{gensymb}
\usepackage{comment}
\usepackage[breaklinks=true]{hyperref}
\usepackage{tkz-euclide} 
\usepackage{listings}
% \usepackage{gvv}                                        
\def\inputGnumericTable{}                                 
\usepackage[latin1]{inputenc}                                
\usepackage{color}                                            
\usepackage{array}                                            
\usepackage{longtable}                                       
\usepackage{calc}                                             
\usepackage{multirow}                                         
\usepackage{hhline}                                           
\usepackage{ifthen}                                           
\usepackage{lscape}
\begin{document}

\bibliographystyle{IEEEtran}

\title{4.13.27}
\author{EE25BTECH11022 - sankeerthan}
% \maketitle
% \newpage
% \bigskip
{\let\newpage\relax\maketitle}

\renewcommand{\thefigure}{\theenumi}
\renewcommand{\thetable}{\theenumi}
\setlength{\intextsep}{10pt} % Space between text and floats


\numberwithin{align}{enumi}
\numberwithin{figure}{enumi}
\renewcommand{\thetable}{\theenumi}

\textbf{Question}: 
Consider the set of all lines $px + qy + r = 0$ such that $3p + 2q + 4r = 0$.
Which one of the following statements is true?
\begin{enumerate}
  \item The lines are concurrent at the point $\brak{\frac{3}{4},\frac{1}{2}}$.
  \item Each line passes through the origin.
  \item The lines are all parallel.
  \item The lines are not concurrent.
\end{enumerate}

\textbf{Solution}:
Given the set of lines $px + qy + r = 0$ such that $3p + 2q + 4r = 0$.
\begin{align}
\myvec{x & y & 1 \\3 & 2 & 4}\myvec{p \\ q \\ r}=\myvec{0 \\ 0}
\end{align}
For a non-trivial solution,the vector $\brak{p, q, r}$ must be in the nullspace. The system has a non-trivial solution if and only if the row vectors are linearly dependent:
\begin{align} 
\brak{x_0, y_0, 1} = \lambda\brak{3, 2, 4} \\
x_0 = 3\lambda,y_0 = 2\lambda,1 = 4\lambda  \\
1 = 4\lambda \implies \lambda = \frac{1}{4} \\
x_0 = \frac{3}{4},\qquad y_0 = \frac{1}{2} 
\end{align}
Therefore, all lines in the given family are concurrent at the point
$\brak{\frac{3}{4},\frac{1}{2}}$

\begin{figure}
    \centering
    \includegraphics[width=1\columnwidth]{figs/points.png}
    \caption{}
    \label{}
\end{figure}
\end{document}